\documentclass[webpdf,contemporary,large,namedate]{oup-authoring-template}%

%\PassOptionsToPackage{hyphens}{url}
%\PassOptionsToPackage{colorlinks,linkcolor=blue,urlcolor=blue,citecolor=blue,anchorcolor=blue}{hyperref}

\DeclareMathOperator*{\argmax}{argmax}

\usepackage{algorithmicx}
\usepackage{lmodern}
\usepackage{setspace}
\renewcommand{\ttdefault}{cmtt}

\begin{document}
\journaltitle{TBD}
\DOI{TBD}
\copyrightyear{2026}
\pubyear{2026}
\access{Advance Access Publication Date: Day Month Year}
\appnotes{Preprint}
\firstpage{1}

\title[Fast read alignment with minibwa]{Fast genomic read alignment with minibwa}
\author[1,2,3,$\ast$]{Heng Li\ORCID{0000-0003-4874-2874}}
\author[4]{Nils Homer\ORCID{0009-0007-7860-1155}}
\address[1]{Department of Data Science, Dana-Farber Cancer Institute, 450 Brookline Ave, Boston, MA 02215, USA}
\address[2]{Department of Biomedical Informatics, Harvard Medical School, 10 Shattuck St, Boston, MA 02215, USA}
\address[3]{Broad Insitute of MIT and Harvard, 415 Main St, Cambridge, MA 02142, USA}
\address[4]{Fulcrum Genomics LLC, 240 Elm Street, Somerville, MA, 02144, USA}
\corresp[$\ast$]{Corresponding author. \href{mailto:hli@ds.dfci.harvard.edu}{hli@ds.dfci.harvard.edu}}

%\received{Date}{0}{Year}
%\revised{Date}{0}{Year}
%\accepted{Date}{0}{Year}

\abstract{
\sffamily\footnotesize
\textbf{Motivation:}
BWA-MEM remains a popular short-read mapper especially for the purpose of variant calling.
Several groups have accelerated this algorithm as it has been the performance bottleneck of many current workflows.
However, constrained by the original design,
these drop-in replacements could only achieve limited speedup.
Breaking changes to BWA-MEM are required for further improvement.\vspace{0.5em}\\
\textbf{Results:}
We developed minibwa for aligning short and accurate long reads against a reference genome.
It combines BWA-MEM variable-length seeding with minimap2 chaining and base alignment.
It speeds up BWA-MEM2 further with additional prefetch for seeding,
new heuristics to skip unnecessary mate rescue and reduced effort in highly repetitive regions
where reads would anyway be wrongly mapped due to structural changes.
Minibwa is over three times as fast as BWA-MEM and over twice as fast as BWA-MEM2 at comparable accuracy.
It also natively supports directional bisulfite sequencing data.
\vspace{0.5em}\\
\textbf{Availability and implementation:}
\url{https://github.com/lh3/minibwa}
}

\maketitle

\section{Introduction}

Read alignment is indispensable to most reference-based sequence analyses such as variant calling and methylation profiling.
Although over 100 read mappers have been published~\citep{Alser:2021aa,DBLP:journals/csur/ProusalisGKNKAPSG26},
only several are maintained and remain widely used.
The two most notable short-read mappers for DNA sequencing data are Bowtie2~\citep{Langmead:2012fk} and BWA-MEM~\citep{Li:2013aa}, both based on FM-index~\citep{DBLP:conf/focs/FerraginaM00}.
While Bowtie2 is as popular for general purposes, BWA-MEM is more often used for variant calling.

BWA-MEM is the sole performance and cost bottleneck in variant calling pipelines.
Given that human Whole-Genome Sequencing (WGS) accounts for the vast majority of sequencing data volume, this inefficiency is particularly pronounced.
There have been several efforts to accelerate the original BWA-MEM algorithm.
BWA-MEM2~\citep{DBLP:conf/ipps/VasimuddinMLA19} revamps the memory layout of FM-index
and accelerates extension alignment with Single Instruction Multiple Data (SIMD) based algorithm.
It is typically 50--100\% faster and outputs identical alignment.
BWA-MEME~\citep{Jung:2022aa} further speeds up BWA-MEM2 with learned index and still promises identical output.
Sentieon BWA is a proprietary drop-in replacement of BWA-MEM.
Parabricks~\citep{OConnell:2023aa,Zhu2025.07.23.666378} provides a GPU-based partial reimplementation of BWA-MEM
that produces alignment functionally equivalent~\citep{Regier:2018aa} to BWA-MEM alignment.

However, designed more than a decade ago, BWA-MEM has several limitations that prevent further optimization but can be alleviated if we accept breaking changes.
First, BWA-MEM uses a complex algorithm to seed alignment~\citep{Li:2012fk}.
It is challenging to apply the latency hiding technique, which we will explain later, to speed up this step.
Second, BWA-MEM chaining is heuristic and inferior to the minimap2~\citep{Li:2018ab} chaining algorithm.
It likely does not work well for short reads longer than 500bp, which will become more common with Roche's SBX sequencing technology~\citep{Kokoris2025.02.19.639056}.
The boundary between short and long reads will be blurred in future.
Third, BWA-MEM uses a non-SIMD-based extension algorithm with intricate boundary conditions.
This greatly complicates SIMD-based equivalence in BWA-MEM2 and increases maintenance burden.
Fourth, BWA-MEM tries too hard in highly repetitive regions.
Due to rapid evolution in centromeres~\citep{Gao2025.12.09.693231}, 
exact base alignment in centromeres is challenging even with complete assemblies~\citep{Bzikadze:2022aa}.
Exhaustive alignment in centromeres is a waste of effort
as we cannot accurately place centromeric short reads anyway due to large structural differences.

In addition to algorithms for aligning standard DNA sequencing data,
there are also many algorithms for aligning bisulfite sequencing (BS-seq) reads.
These include
BSMAP~\citep{Xi:2009aa},
Bismark~\citep{Krueger:2011aa},
BS-Seeker~\citep{Chen:2010aa,Guo:2013aa,Huang:2018aa},
BWA-Meth~\citep{Pedersen:2014aa},
BitMapperBS~\citep{Cheng:2018aa},
gemBS~\citep{Merkel:2019aa},
HISAT-3N~\citep{Zhang:2021ac},
BSBolt~\citep{Farrell:2021aa},
abismal~\citep{de_Sena_Brandine_2021},
FAME~\citep{Fischer_2023},
BISCUIT~\citep{Zhou:2024ac},
RMapAlign3N~\citep{Muller:2025aa}
and ARYANA-BS~\citep{Nikaein_2025}.
Several generic short-read aligners, such as
GSNAP~\citep{Wu:2010uq},
segemehl~\citep{Otto:2014aa}
and Arioc~\citep{Wilton:2018aa},
also natively support BS-seq.

Among these algorithms, Bismark~\citep{Krueger:2011aa} is by far the most cited,
partly because it is a wrapper around stock Bowtie~\citep{Langmead:2009aa} and Bowtie2~\citep{Langmead:2012fk}
which are battle tested and whose behavior is well understood.
BWA-Meth~\citep{Pedersen:2014aa} is a popular wrapper around BWA-MEM
and has been shown to have good performance in recent benchmarks~\citep{Foox:2021aa,Gong_2022,Kerns_2025,Lin_2025}.
It concatenates C-to-T and G-to-A converted reference genomes and builds one FM-index.
For paired-end reads produced with the directional protocol,
BWA-Meth converts {\tt C} bases on read 1 to {\tt T} and converts {\tt G} bases on read 2 to {\tt A}.
It aligns converted reads to the FM-index with stock BWA-MEM
and postprocesses the BWA-MEM output to generate the final alignment.
BWA-Meth has several limitations as a wrapper.
First, with the directional protocol, read 1 should be mapped either to the forward strand of the C-to-T reference or the reverse strand of the G-to-A reference.
Unaware of this constraint, BWA-MEM will also attempt to map read 1 to the reverse C-to-T and forward G-to-A reference genomes.
Although these hits are later filtered out by BWA-Meth,
they waste computing time and may affect mapping quality calculation.
Second, if a read 1 is mapped to the C-to-T reference, read 2 will be expected to get mapped to the G-to-A reference.
The two reads in a pair are not locate on the same contig.
This breaks BWA-MEM read pairing.
Third, BWA-MEM attempts to rescue a read if its mate is well aligned.
However, because the two reads are not supposed to be aligned to the same contig due to base conversion,
mate rescue would often fail.

BISCUIT~\citep{Zhou:2024ac} and BSBolt~\citep{Farrell:2021aa} are also built upon BWA-MEM.
They avoid the issues mentioned above by customizing BWA-MEM source code specifically for BS-seq data.
They are theoretically more rigorous than BWA-Meth.
As a tradeoff, though, the deep code customization makes it challenging to adopt the faster BWA-MEM2 implementation.
Combining an efficient mapping algorithm with the BISCUIT/BSBolt strategy will yield an even better BS-seq mapper.

\section*{Acknowledgments}

blabla

\section*{Author contributions}

H.L. conceived the project, implemented the core algorithm, analyzed the data and drafted the manuscript.
N.H. added new features, contributed to development and reviewed the draft.

\section*{Conflict of interest}

N.H. is an employee of Fulcrum Genomics LLC.

\section*{Funding}

This work is supported by National Institute of Health grant R01HG010040, R01HG014175 and U24CA294203 (to H.L.).

\section*{Data availability}

The minibwa source code is available at \url{https://github.com/lh3/minibwa}.

\bibliographystyle{apalike}
{\sffamily\small
\bibliography{minibwa}}

\end{document}
